%================================================================
\appendix

\section{NormCode Syntax Reference}
\label{app:syntax}

This appendix provides a concise reference for the NormCode grammar used throughout the
paper. The complete formal specification is in \cite{anon2025}.

\subsection*{A.1 Concept Markers}

\begin{table}[h]
\centering
\caption{The three NormCode concept markers.}
\label{tab:markers}
\begin{tabular}{@{}p{1.4cm}p{2.4cm}p{4.0cm}p{3.4cm}@{}}
\toprule
\textbf{Marker} & \textbf{Type} & \textbf{Meaning} & \textbf{CBR role} \\
\midrule
\texttt{<-} & Value Concept   & A named value, document, or intermediate result --- data flowing between steps & Stored in Concept Repository; constitutes $\mathcal{C}_f$ \\[4pt]
\texttt{<=} & Functional Concept & An operation deriving its parent concept from declared inputs; semantic (LLM) or syntactic (deterministic) & Produces checkpoints; maps to flow indices \\[4pt]
\texttt{<*} & Context Concept & Current iteration element; carries in-loop state without leaking to outer scopes & Scopes loop-local state; prevents cross-iteration contamination \\
\bottomrule
\end{tabular}
\end{table}

\subsection*{A.2 Syntactic Operator Families}

Syntactic inferences (\texttt{<=} with a syntactic sequence type) belong to one of four
families. All are deterministic and zero LLM cost; their checkpoints are perfectly
reproducible Level~1 cases.

\begin{table}[h]
\centering
\caption{Syntactic operator families.}
\label{tab:operators}
\begin{tabular}{@{}p{1.8cm}p{1.2cm}p{4.4cm}p{4.0cm}@{}}
\toprule
\textbf{Family} & \textbf{Symbol} & \textbf{Function} & \textbf{Example use} \\
\midrule
Assigning & \$  & Routes or binds a value concept to another concept & Pass a ground input into a nested scope \\[4pt]
Grouping  & \&  & Collects or concatenates a set of concepts into a single reference & Collect loop outputs into \texttt{all slides} \\[4pt]
Timing    & @   & Gates execution on a condition or sequencing dependency & Wait for parallel branches before proceeding \\[4pt]
Looping   & *   & Manages iteration over a collection with carried \texttt{<*} state & Per-document summarization; per-slide generation \\
\bottomrule
\end{tabular}
\end{table}

\subsection*{A.3 Flow Index Assignment (Sibling Pattern)}

Flow indices are assigned during the Formalization phase. The rule is:
\begin{itemize}
  \item The root output concept receives index \texttt{1}.
  \item The functional concept (\texttt{<=}) under any parent is always \texttt{.1}.
  \item Value inputs (\texttt{<-}) are siblings at \texttt{.2}, \texttt{.3}, \ldots in
    declaration order.
  \item Nesting is recursive: index \texttt{1.2.3} denotes the third declared input of
    the second declared input of the root.
\end{itemize}

For the two-concept example in Sect.~\ref{sec:normcode}:
\texttt{summary} $\rightarrow$ \texttt{1},\quad
\texttt{<=\ summarize} $\rightarrow$ \texttt{1.1},\quad
\texttt{report} $\rightarrow$ \texttt{1.2},\quad
\texttt{<=\ read} $\rightarrow$ \texttt{1.2.1},\quad
\texttt{source\_doc} $\rightarrow$ \texttt{1.2.2},\quad
\texttt{style\_guide} $\rightarrow$ \texttt{1.3}.

Every checkpoint in $\mathcal{CB}$ is keyed by (\texttt{run\_id},
\texttt{flow\_index}), making any past runtime state directly addressable.

\subsection*{A.4 \texttt{.ncn} Translation Tags}

The Formalization phase generates a \texttt{.ncn} plain-prose narrative using four
translation tags:

\begin{table}[h]
\centering
\caption{\texttt{.ncn} translation tags.}
\label{tab:ncntags}
\begin{tabular}{@{}p{2.2cm}p{9.0cm}@{}}
\toprule
\textbf{Tag} & \textbf{Meaning in the generated narrative} \\
\midrule
\texttt{(OUTPUT)}  & Root output concept --- the final product of the plan \\
\texttt{(ACTION)}  & A functional concept --- what is being done at this step \\
\texttt{(VALUE)}   & A value concept --- a named data item or intermediate result \\
\texttt{(CONTEXT)} & A context concept --- the current element in a loop iteration \\
\bottomrule
\end{tabular}
\end{table}

%================================================================
\section{Base-X Addition Benchmark Plan}
\label{app:basex}

This appendix shows the NormCode plan used for the base-X addition benchmark
(Sect.~\ref{sec:evaluation}). The plan adds two arbitrary-base numbers digit-by-digit
with carry-over, structured as three nested loops. It validates both execution
correctness (100\% accuracy up to 150-digit numbers) and Level~1 CBR (checkpoints at
each loop iteration are inspectable and revisable via failure localization (C1) and selective re-execution (C3)).

\begin{lstlisting}[caption={Base-X addition plan (\texttt{.ncds} format).},
                   label={lst:basex}]
<- addition result
    <= report final sum in target base
    <- digit sums
        <= for each number pair in the collection
        <- digit sum
            <= add digits with carry-over
            <- digit pair
                <= extract unit-place digits
                <- shifted pair
                    <= generate shifted number pair
                    <- number pair
                <- base
            <- carry
        <- number pairs
        <* number pair
    <- number pairs
        <= for each input number
        <- number pair
            <= pair numbers at current position
            <- number_a
            <- number_b
        <- numbers
        <* current_number
    <- numbers
    <- base
\end{lstlisting}

The three nested loops correspond to: (1)~outer loop over digit positions across all
number pairs; (2)~inner loop extracting unit-place digits from shifted pairs; (3)~carry
propagation loop resolving digit-by-digit addition. The \texttt{base} concept is
declared as a ground input at the outermost scope and passed explicitly into each nested
scope --- the scope rule prevents any loop body from accessing base implicitly through
accumulated state.

\textbf{Level~1 validation.} Mid-computation checkpoints can be inspected at any loop
iteration: the Tensor Inspector shows the current digit, carry value, and accumulated
sum at that flow index. Overriding a carry value at flow index \texttt{1.1.3.2} (for
example) and resuming triggers selective re-execution of only the downstream digit
addition steps --- selective re-execution (C3) with exactly the minimum necessary re-execution.
